\chapter{Gaussian tails and convergence of entropy}\label{sec:entropy-convergence}
Local convergence of measures controls compact tests. Gaussian integrals need tail control, and entropy adds a supremum over a noncompact parameter space. We separate these steps. Use $F_{c,t}$ and $\lambda$ from Definition \ref{def:mcf-entropy}, with dimension $n\geq1$, and write $\rho_{c,t}(x)=(4\pi t)^{-n/2}e^{-|x-c|^2/(4t)}$. The area-growth estimates extend the argument of \cite{CM12}, Lemma 7.28, pp.~788--789, to Radon measures.

\begin{proposition}[Lower semicontinuity of Gaussian area and entropy]\label{prop:entropy-lower-semicontinuity}
If nonnegative Radon measures $\mu_j$ on $\mathbb R^{n+1}$ converge locally weakly to $\mu$, then for every $c$ and $t>0$,
\[
 F_{c,t}(\mu)\leq\liminf_j F_{c,t}(\mu_j),\qquad
 \lambda(\mu)\leq\liminf_j\lambda(\mu_j).
\]
The values may be infinite, and no uniform growth bound is required.
\end{proposition}
\begin{proof}
Choose compactly supported continuous cutoffs $0\leq\chi_R\leq1$ increasing to one. Local weak convergence gives $\int\chi_R\rho_{c,t}\,d\mu=\lim_j\int\chi_R\rho_{c,t}\,d\mu_j\leq\liminf_jF_{c,t}(\mu_j)$. Monotone convergence in $R$ proves the first inequality. The same argument bounds each $F_{c,t}(\mu)$ by $\liminf_j\lambda(\mu_j)$; take the supremum in $c,t$. This last step does not interchange a limit and a supremum as an equality.
\end{proof}

\begin{lemma}[Uniform polynomial Gaussian tails]\label{lem:entropy-gaussian-tails}
Suppose $\lambda(\mu)\leq\Lambda$. Fix $q\geq0$, $A\geq0$ and $0<a\leq b<\infty$. Then
\[
 \sup_{|c|\leq A,\ a\leq t\leq b}
 \int_{|x|>R}(1+|x|)^q\rho_{c,t}(x)\,d\mu(x)\longrightarrow0
\]
as $R\to\infty$, uniformly over all such measures. More precisely, for $R\geq2A+2$ the integral is at most
\[
 C(n,a)\Lambda\sum_{j=\lfloor R\rfloor}^{\infty}
 (j+2)^{n+q}\exp\left(-\frac{(j-A)^2}{4b}\right).
\]
\end{lemma}
\begin{proof}
Proposition \ref{prop:mcf-entropy-area-growth} bounds $\mu(B_{j+1}(0))$ by $C(n)\Lambda(j+1)^n$. In the annulus $j\leq|x|<j+1$, the polynomial factor is at most $(j+2)^q$, $|x-c|\geq j-A$, and the prefactor is at most $(4\pi a)^{-n/2}$. Sum these bounds. The series converges, with tails tending to zero independently of the measure and parameters. The same estimate with the bounded inner ball included gives uniform finite polynomial Gaussian moments. This is the annular method of \cite{CM12}, (7.31)--(7.32).
\end{proof}

\begin{theorem}[Gaussian convergence on compact parameter sets]\label{thm:gaussian-local-weak-convergence}
Suppose nonnegative Radon measures $\mu_j\rightharpoonup\mu$ locally in $\mathbb R^{n+1}$ and $\lambda(\mu_j)\leq\Lambda$. Then
\[
 F_{c,t}(\mu_j)\longrightarrow F_{c,t}(\mu)
\]
uniformly for $(c,t)$ in every compact subset of $\mathbb R^{n+1}\times(0,\infty)$. In particular the conclusion applies to the weight measures of locally convergent varifolds with uniformly bounded entropy. This uses the Gaussian tail estimates associated with \cite{CM12}, Lemma 7.28.
\end{theorem}
\begin{proof}
Proposition \ref{prop:entropy-lower-semicontinuity} gives $\lambda(\mu)\leq\Lambda$. For fixed $(c,t)$, insert a cutoff on $B_R$; local weak convergence handles the cutoff Gaussian, while Lemma \ref{lem:entropy-gaussian-tails} makes both omitted tails uniformly small. This proves pointwise convergence.

On $|c|\leq A$, $a\leq t\leq b$, the parameter derivatives of the kernel are
\[
 \nabla_c\rho_{c,t}=\frac{x-c}{2t}\rho_{c,t},\qquad
 \partial_t\rho_{c,t}=\left(\frac{|x-c|^2}{4t^2}-\frac n{2t}\right)\rho_{c,t}.
\]
The moment bounds give a uniform bound for their integrals on a slightly larger compact parameter box. Differentiation under the integral follows by the same domination. Thus all $F_{c,t}(\mu_j)$ and the limit function have a common Lipschitz constant there. Approximate the compact parameter set by a finite net: pointwise convergence on the net and the common Lipschitz bound give uniform convergence. No control of parameters with $t\to0$, $t\to\infty$, or $|c|\to\infty$ is asserted.
\end{proof}

\begin{example}[Entropy can drop when a plane escapes to infinity]\label{ex:entropy-two-plane-escape}
Let $P_h=\{x_{n+1}=h\}$ and $\mu_j=\mathcal H^n\llcorner P_0+\mathcal H^n\llcorner P_j$. Then $\mu_j\rightharpoonup\mathcal H^n\llcorner P_0$ locally, while $\lambda(\mu_j)=2$ and $\lambda(P_0)=1$. Thus entropy is not upper semicontinuous even with a common entropy bound and a fixed component.
\end{example}
\begin{proof}
The second plane eventually misses every compact set. If the center has last coordinate $a$, Example \ref{ex:mcf-plane-entropy} gives
\[
 F_{c,t}(\mu_j)=e^{-a^2/(4t)}+e^{-(a-j)^2/(4t)}\leq2.
\]
For fixed $j,c$ this tends to two as $t\to\infty$, proving the entropy value. At each fixed finite center and scale, the second term tends to zero as $j\to\infty$, in agreement with Theorem \ref{thm:gaussian-local-weak-convergence}. The optimizing scales are not confined to a compact parameter set.
\end{proof}

\begin{proposition}[Entropy controls the Gaussian second moment]\label{prop:entropy-gaussian-second-moment}
For every nonnegative Radon measure with finite entropy,
\[
 \int |x|^2\rho_{c,t}(x)\,d\mu(x)
 \leq2\lambda(\mu)\bigl(|c|^2+C_n t\bigr),\qquad t>0.
\]
Here $C_n$ depends only on the dimension. This is the measure version of \cite{CM12}, Lemma 7.28.
\end{proposition}
\begin{proof}
Apply the geometric translation and dilation $z=(x-c)/\sqrt t$ from Proposition \ref{prop:mcf-entropy-invariance}, including the factor $t^{-n/2}$ in the transformed measure $\widetilde\mu$. The integral becomes $\int|c+\sqrt t z|^2\rho_{0,1}(z)\,d\widetilde\mu$. The inequality $|c+\sqrt t z|^2\leq2|c|^2+2t|z|^2$ bounds its first part by $2|c|^2\lambda(\mu)$. The area bound and unit annular sum in Lemma \ref{lem:entropy-gaussian-tails} bound the second moment of $\widetilde\mu$ by $C_n\lambda(\widetilde\mu)=C_n\lambda(\mu)$.
\end{proof}

\begin{proposition}[Entropy is unchanged by a Euclidean product]\label{prop:entropy-euclidean-product}
Let $\mu$ be a nonnegative Radon measure on $\mathbb R^{m+1}$ with Gaussian normalization of dimension $m\geq1$. For $k\geq1$, give $\mu\times\mathcal L^k$ the normalization of dimension $m+k$. Then
\[
 F^{m+k}_{(c,b),t}(\mu\times\mathcal L^k)=F^m_{c,t}(\mu),
 \qquad \lambda_{m+k}(\mu\times\mathcal L^k)=\lambda_m(\mu).
\]
\end{proposition}
\begin{proof}
The nonnegative product Gaussian factors. Tonelli's theorem and $\int_{\mathbb R^k}(4\pi t)^{-k/2}e^{-|z-b|^2/(4t)}\,dz=1$ give the first equality, even when the other factor is infinite. Taking suprema gives the second. Translation of the Gaussian center in a product direction has no effect; this will matter in Proposition \ref{prop:shrinker-gaussian-equality}.
\end{proof}

\begin{proposition}[Small-scale Gaussian limits on a proper embedded hypersurface]\label{prop:static-gaussian-density}
Let $\Sigma^n\subset\mathbb R^{n+1}$ be smooth, properly embedded, without boundary, and with polynomial area growth. Then
\[
 \lim_{t\downarrow0}F_{c,t}(\Sigma)=
 \begin{cases}1,&c\in\Sigma,\\0,&c\notin\Sigma.\end{cases}
\]
In particular $\lambda(\Sigma)\geq1$ when $\Sigma$ is nonempty. This is the local statement in \cite{CM12}, Lemma 7.2(3), with properness explicit.
\end{proposition}
\begin{proof}
At $c\in\Sigma$, represent a neighborhood as one smooth graph over its tangent plane, with graph value and derivative zero at $c$. Rescaling by $t^{-1/2}$ gives convergence to that plane on bounded sets. In a fixed small graph neighborhood the area Jacobian is bounded and the tangential Gaussian dominates an integrable Gaussian on the tangent plane; hence its integral tends to one by dominated convergence. Proper embeddedness allows this neighborhood to exclude all other sheets.

Outside a fixed ambient ball of radius $\delta$ about $c$, for $0<t\leq t_*$ the Gaussian contribution is bounded by
\[
 (4\pi t)^{-n/2}e^{-\delta^2/(8t)}
 \int_\Sigma e^{-|x-c|^2/(8t_*)}\,dA,
\]
which tends to zero; the last integral is finite by polynomial area growth. If $c\notin\Sigma$, properness makes $\Sigma$ closed, so its distance from $c$ is positive and the same estimate applies to the whole surface. This is the static counterpart of Proposition \ref{prop:mcf-regular-density}; multiplicity or nonproper accumulation would change the argument.
\end{proof}
